\documentclass[a4paper,11pt]{article}
\usepackage[utf8]{inputenc}
\usepackage[T1]{fontenc}
\usepackage[french]{babel}
\usepackage[left=2cm,right=2cm,top=2cm,bottom=2cm]{geometry}
\usepackage{xcolor}
\usepackage{hyperref}
\usepackage{fontawesome5}
\usepackage{parskip}

% --- Couleurs Servier ---
\definecolor{primary}{RGB}{52, 72, 138}
\hypersetup{colorlinks=true, urlcolor=primary}

\begin{document}

% --- EXPÉDITEUR ---
\noindent
\begin{minipage}[t]{0.5\textwidth}
    \textbf{Loïc Aron MBASSI EWOLO}\\
    Élève-Ingénieur Bac+5 -- Double diplôme ENSEA\\[3pt]
    \faMapMarker\ Cergy, France\\
    \faPhone\ (+33) 7 59 52 33 98\\
    \faEnvelope\ \href{mailto:loicaron.mbassiewolo@ensea.fr}{loicaron.mbassiewolo@ensea.fr}
\end{minipage}
\hfill
\begin{minipage}[t]{0.4\textwidth}
    \raggedleft
    Cergy, le 20 mars 2026
\end{minipage}

\vspace{1.2cm}

\hfill
\begin{minipage}[t]{0.5\textwidth}
    \raggedleft
    \textbf{Servier}\\
    Service Recrutement\\
    Suresnes, France
\end{minipage}

\vspace{1cm}

\noindent\textbf{Objet :} Candidature au poste d'Alternant Data Engineer (H/F) -- Contrat de Professionnalisation -- N\textdegree{} offre 11222

\vspace{0.8cm}

Madame, Monsieur,

Concevoir des pipelines de données Python à grande échelle, analyser des politiques de confidentialité par NLP, modéliser des scores interprétables pour des chercheurs et des utilisateurs finaux : c'est ce que je fais actuellement à l'INRIA Saclay, au sein de l'Institut Polytechnique de Paris. C'est aussi ce qui me conduit à vous adresser ma candidature pour le contrat de professionnalisation Data Engineer au sein de votre équipe Data Engineering.

Mon stage à l'INRIA Saclay (Équipe TRIBE, programme national PEPR MOBIDEC, Avril-Août 2026) m'a confronté à des enjeux concrets de traitement de données à grande échelle : conception d'un pipeline d'analyse automatisée combinant NLP (BERT/Transformers), analyse statique de code et modélisation de scores, avec pandas et scikit-learn comme socle technique. Cette expérience m'a appris à travailler sur des données réelles, hétérogènes et volumineuses, dans un environnement de recherche exigeant en termes de qualité et de reproductibilité. Avant cela, mon rôle d'Assistant de Recherche au MILA (Institut Québécois d'IA, été 2025) m'a formé à l'analyse statistique rigoureuse et à la visualisation des dynamiques d'entraînement de modèles à grande échelle.

Sur le plan applicatif, la conception d'un système multi-agents agricole (Google ADK, LangChain, FastAPI, Docker, orchestration d'APIs externes dont des services Google Cloud) m'a permis de construire et de déployer un pipeline de données complet, de l'ingestion à la restitution. Mon projet UROP en analyse d'ECG, mené en collaboration avec des mentors Google DeepMind, m'a confronté aux exigences d'un pipeline médical : qualité des données, extraction de features sur séries temporelles, contraintes de reproductibilité. Ces deux expériences illustrent ma capacité à passer de la donnée brute à une valeur opérationnelle dans des contextes métier exigeants, ce qui correspond précisément aux missions décrites dans votre offre.

Je suis disponible dès septembre 2026, à l'issue de mon stage INRIA, pour rejoindre votre équipe Data Engineering à Suresnes dans le cadre d'un contrat de professionnalisation. Servier représente un contexte particulièrement stimulant : la donnée y sert directement le progrès thérapeutique, ce qui donne un sens concret à chaque pipeline construit. Je serais ravi d'échanger avec vous pour vous exposer plus en détail ma motivation et mes compétences.

Je vous prie d'agréer, Madame, Monsieur, l'expression de mes salutations distinguées.

\vspace{0.8cm}

\textbf{Loïc Aron MBASSI EWOLO}\\
\textit{Élève-Ingénieur ENSEA / Polytechnique Yaoundé}

\end{document}
